\documentclass{article}
\usepackage[T1]{fontenc}
\usepackage{lmodern}
\usepackage{graphicx}
\usepackage{amsmath,amssymb}
\usepackage[a4paper,margin=2cm]{geometry}
\usepackage[absolute,overlay]{textpos}
\setlength{\TPHorizModule}{1mm}
\setlength{\TPVertModule}{1mm}
\usepackage[indonesian]{babel}
\usepackage{hyperref}
\setlength{\emergencystretch}{2em}
% unit: Halaman 70

\begin{document}

% === Halaman 70 (python/native) ===

Pembangkitan kunci putaran

• Setiap putaran di dalam iterated cipher menggunakan kunci internal (subkey) 

yang dinamakan kunci putaran (round key).

• Kunci putaran dibangkitkan dari kunci eksternal yang diberikan oleh user melalui 

proses yang dinamakan key expansion atau key scheduling.

• Di dalam key expansion dilakukan komputasi yang kompleks untuk menghasilkan 

sejumlah kunci putaran yang berbeda-beda. 

• Panjang kunci putaran tidak selalu sama dengan panjang kunci eksternal.  

70

\end{document}
